La sévérité de la dépression se mesure par le score PHQ-9, et des biomarqueurs sanguins pourraient aider à la suivre. Nous reproduisons le protocole de Bader \emph{et al.} (2024), fondé sur une forêt aléatoire pondérée (WRF), sur 2\,621 observations décrites par trois biomarqueurs mitochondriaux et neuf variables cliniques, puis nous l'améliorons par un pipeline complet d'optimisation d'hyperparamètres (GridSearchCV, RandomizedSearchCV, recherche bayésienne Optuna) sous validation croisée stratifiée. Dix modèles sont comparés avant et après optimisation, en classification binaire et à quatre classes, avec une validation croisée imbriquée $5\times 2$ et des tests de Wilcoxon, de Nadeau--Bengio et de Friedman--Conover. Dans le protocole de l'article, la RF atteint une AUC de test de 0,93 (binaire) et 0,90 (quatre classes), et la WRF optimisée 0,93 ; mais l'optimisation ne profite significativement qu'à KNN, SVM et ANN ($+0{,}04$ à $+0{,}12$ d'AUC), pas aux forêts. Un audit montre que 67 à 85\,\% des valeurs de biomarqueurs sont dupliquées entre lignes et presque parfaitement alignées sur la classe de sévérité : sur les 228 lignes sans doublon, l'AUC chute à 0,60--0,72 et l'optimisation ne procure plus aucun gain (RF : $-0{,}03$). Les scores élevés semblent donc refléter une fuite de données plutôt qu'un signal biologique, ce qui conditionne toute conclusion.
